\section{Related Work}

\paragraph{Truncation and timeout semantics.}
Time limits used to facilitate training should not be conflated with terminal states of the underlying Markov decision process. Pardo et al.\ formalize this distinction and show that partial episodes require value bootstrapping at truncation \cite{pardo2018timelimits}. Contemporary environment APIs distinguish termination from truncation for the same reason, and D4RL similarly exposed timeout metadata in offline datasets \cite{fu2020d4rl}. CVT applies this known principle to trajectory-aware multi-step offline GCRL; the novelty of SegBench-GC is the controlled invariance benchmark, not the Bellman correction itself.

\paragraph{Offline goal-conditioned RL.}
OGBench evaluates offline GCRL under stitching, long-horizon reasoning, mixed-quality data, and high-dimensional observations \cite{ogbench2024}. Its datasets distinguish trajectory ends from Bellman masks: trajectory boundaries constrain future-goal sampling, while masks indicate task success. SegBench-GC preserves that distinction and adds a separate set of artificial backup boundaries, allowing target sensitivity to be measured without changing the goal distribution.

\paragraph{Temporal abstraction and multi-step value propagation.}
Implicit Q-learning and conservative Q-learning learn from fixed datasets without online exploration \cite{kostrikov2022iql,kumar2020cql}. More recent offline-GCRL work directly targets long effective horizons. Horizon Reduction studies several ways to shorten the learning horizon, including n-step GCSAC+BC \cite{park2025horizon}; OTA incorporates option-aware temporal abstraction into value learning \cite{ahn2025ota}; multistep quasimetric learning combines multistep propagation with quasimetric structure \cite{zheng2026mqe}; Decoupled Q-Chunking separates critic and policy chunk lengths while retaining multi-step value propagation \cite{li2026dqc}; and Transitive RL replaces linear-depth TD propagation with a divide-and-conquer value update \cite{park2026trl}. In their reported evaluations, these methods do not use benign resegmentation as a counterfactual reliability intervention while holding transitions and future-goal semantics fixed.

\paragraph{Segmentation and information loss.}
Recent theory shows that restricting offline RL to fixed-length trajectory segments can cause identifiability failure, objective misspecification, and support aliasing \cite{nidadala2026horizon}. SegBench-GC studies a complementary setting: ordered transitions remain available, source trajectories remain intact for goal sampling, and only a separate backup segmentation is changed. Hierarchical action-chunking work likewise uses $k$-step value backups for long-horizon offline GCRL \cite{jawaid2026hiqc}, but does not isolate whether an administrative backup cut retains continuation value. We therefore treat SegBench-GC as a diagnostic protocol for learners that consume trajectory prefixes while limiting empirical claims to the two tested learners.

\paragraph{Reliable empirical evaluation.}
Deep RL comparisons are highly variable and should report paired seeds, uncertainty, and distributions rather than isolated point estimates \cite{agarwal2021statistical}. We pair optimization seeds and segmentation configurations where possible, aggregate segmentation variants within seed before aggregating across seeds, and report paired contrasts alongside descriptive variability.
